\documentclass[aspectratio=169,11pt]{beamer}

\usepackage{fontspec}
\usetheme{metropolis}
\usepackage{booktabs}

% ---- Fira Sans (explicit paths; not in fontconfig) ----------------------
\setsansfont{Fira Sans}[
  Path=/usr/share/texlive/texmf-dist/fonts/opentype/public/fira/,
  Extension=.otf,
  UprightFont=FiraSans-Regular,
  BoldFont=FiraSans-Bold,
  ItalicFont=FiraSans-Italic,
  BoldItalicFont=FiraSans-BoldItalic,
]
\setmonofont{Fira Mono}[
  Path=/usr/share/texlive/texmf-dist/fonts/opentype/public/fira/,
  Extension=.otf,
  UprightFont=FiraMono-Regular,
  BoldFont=FiraMono-Bold,
]
\renewcommand{\familydefault}{\sfdefault}

% ---- Polkadot pink accent ----------------------------------------------
\definecolor{polkadotpink}{HTML}{E6007A}
\definecolor{ink}{HTML}{1A1A1A}

\setbeamercolor{palette primary}{fg=polkadotpink}
\setbeamercolor{frametitle}{fg=polkadotpink}
\setbeamercolor{progress bar}{fg=polkadotpink}
\setbeamercolor{alerted text}{fg=polkadotpink}
\setbeamercolor{title separator}{fg=polkadotpink}
\setbeamercolor{normal text}{fg=ink}
\metroset{block=fill,progressbar=frametitle,numbering=fraction}

% Keep the title frame within the 16:9 text height
\setbeamerfont{title}{size=\LARGE,series=\bfseries}
\setbeamerfont{subtitle}{size=\large}
\setbeamerfont{author}{size=\normalsize}
\setbeamerfont{institute}{size=\small}
\setbeamerfont{date}{size=\small}

\title{Blockchain in Practice}
\subtitle{Two Video Overviews: Polkadot and Kubernetes}
\author{Aidas Kri\v{s}\v{c}i\={u}nas}
\institute{Blockchain Course}
\date{2026}

\begin{document}

\maketitle

% ---------------------------------------------------------------------------
\begin{frame}{Overview}
  \small
  \begin{itemize}
    \item Two videos from the course material are summarised: content and key concepts only (no rating).
    \item \alert{Video 1} --- from the course playlist:
          \emph{What is Polkadot? A Polkadot for Beginners Guide and Intro to Blockchain}
          (entry 24).
    \item \alert{Video 2} --- free choice:
          \emph{Managing and Securing Blockchain Applications on Kubernetes}
          (Haining Zhang \& Yang Yu, VMware; KubeCon + CloudNativeCon China 2018).
    \item Required viewing only, not reviewed: \emph{How the Blockchain is Changing Money and Business}.
  \end{itemize}
  \vfill
  \footnotesize Each video is presented as a short overview plus its main concepts and 2--3 keywords.
\end{frame}

% ===========================================================================
\begin{frame}{Video 1: What It Is About}
  \small
  \begin{columns}[T,onlytextwidth]
    \begin{column}{0.55\textwidth}
      \textbf{Talk at a glance}
      \begin{itemize}
        \item Channel: \emph{Polkadot}; published May 2020; runtime $\approx$76\,min.
        \item A beginner, jargon-free introduction to Polkadot and blockchain.
        \item Presenters:
          \begin{itemize}
            \item \textbf{Bill Laboon} --- Technical Education Lead, Web3 Foundation.
            \item \textbf{Dan Reecer} --- Community \& Growth, Web3 Foundation.
          \end{itemize}
      \end{itemize}
    \end{column}
    \begin{column}{0.42\textwidth}
      \textbf{Roadmap stated up front}
      \begin{enumerate}
        \item How the Internet came about and evolved.
        \item History of blockchain and crypto, and the similarities.
        \item What blockchain is.
        \item Problems of legacy blockchains.
        \item How Polkadot solves them.
        \item What to expect from Polkadot long term.
      \end{enumerate}
    \end{column}
  \end{columns}
\end{frame}

\begin{frame}{Video 1: From the Internet to the Blockchain}
  \scriptsize
  \textbf{Internet history (the parallel)}
  \begin{itemize}
    \item Late 1960s: computers were ``islands unto themselves'' --- data moved by diskette, tape, punch cards.
    \item \textbf{ARPANET}: first network of computing sites (universities CMU, Rutgers, Illinois, plus e.g.\ NSA).
    \item Early file-sharing and email; Queen Elizabeth II sent her first email in 1976.
    \item \textbf{1978}: Robert Kahn's TCP/IP let separate networks share one language --- the birth of the Internet.
  \end{itemize}
  \textbf{What a block is}
  \begin{itemize}
    \item ``Chunks of data \ldots associated with a particular time period \ldots connected in a chain and distributed among many computers.''
    \item No central place where the blockchain lives; it cannot be altered once the computers agree on the next block.
    \item Analogies: every poker hand recorded by all players; a chess move-sheet anyone can replay.
  \end{itemize}
  \textbf{Chain of history:} 1980s digital cash (eCash) $\rightarrow$ 1997 Hashcash (Adam Back) $\rightarrow$ 2008 Bitcoin paper / 2009 network $\rightarrow$ 2015 Ethereum (programmable apps).
\end{frame}

\begin{frame}{Video 1: Key Concepts \& Keywords}
  \scriptsize
  \begin{columns}[T,onlytextwidth]
    \begin{column}{0.5\textwidth}
      \textbf{Decentralized system}
      \begin{itemize}
        \item No central point of storage; many computers maintain one shared chain.
        \item Lets people trust each other without knowing each other; very hard to shut down.
      \end{itemize}
      \textbf{Problems of legacy chains}
      \begin{itemize}
        \item Poor interoperability --- good within one chain, weak across networks.
        \item One size fits all (no customisation); weak governance; hard upgrades.
        \item Limited scalability: Bitcoin $\approx$3--4 TPS vs.\ Visa $>$1{,}700; high fees when busy.
      \end{itemize}
    \end{column}
    \begin{column}{0.47\textwidth}
      \textbf{How Polkadot solves them}
      \begin{itemize}
        \item \alert{Relay chain + parachains}: many custom chains (IoT, finance, gaming) into one ``organism''.
        \item \alert{Shared security}: validators shared across chains; collators feed data to them.
        \item Runtime upgrades by on-chain governance: proposal $\rightarrow$ second $\rightarrow$ referendum $\rightarrow$ majority of stake.
        \item \emph{Substrate} framework to build chains; \emph{Kusama} as a canary network.
      \end{itemize}
    \end{column}
  \end{columns}
  \vfill
  \textbf{Keywords:} \alert{Internet} \;$\cdot$\; \alert{Definition of Block} \;$\cdot$\; \alert{Decentralized System}.
\end{frame}

% ===========================================================================
\begin{frame}{Video 2: What It Is About}
  \small
  \begin{itemize}
    \item \textbf{Talk:} \emph{Managing and Securing Blockchain Applications on Kubernetes}.
    \item \textbf{Venue:} KubeCon + CloudNativeCon China 2018 (Nov 14, 2018).
    \item \textbf{Speakers (VMware China R\&D):}
      \begin{itemize}
        \item \textbf{Haining ``Henry'' Zhang} --- Technical Director; founder of the Harbor project.
        \item \textbf{Yang Yu} --- Software Engineer; Kubernetes CNI plugin / NSX; ex-OpenStack.
      \end{itemize}
    \item \textbf{Premise:} because blockchain is \emph{decentralised}, the cloud is an ideal place to run it, and Kubernetes is a \emph{ubiquitous cloud operating system}.
    \item \textbf{Goal:} reduce operational complexity by automating enterprise blockchain deployment.
  \end{itemize}
\end{frame}

\begin{frame}{Video 2: Key Concepts \& Keywords}
  \footnotesize
  \begin{columns}[T,onlytextwidth]
    \begin{column}{0.5\textwidth}
      \textbf{Enterprise / private blockchain}
      \begin{itemize}
        \item \alert{Hyperledger Fabric} --- an open-source, permissioned blockchain framework.
        \item Aimed at business problems, not public cryptocurrency use.
      \end{itemize}
      \textbf{Automated deployment on Kubernetes}
      \begin{itemize}
        \item Deploy Hyperledger Fabric automatically on Kubernetes.
        \item Kubernetes as the ubiquitous cloud OS reduces operational complexity.
      \end{itemize}
    \end{column}
    \begin{column}{0.47\textwidth}
      \textbf{Security by isolation}
      \begin{itemize}
        \item Use a \alert{CNI plugin} to apply network policies based on Kubernetes \emph{namespaces}.
        \item Isolates workloads between blockchain entities (peers, orderers, organisations).
      \end{itemize}
      \textbf{Outcome}
      \begin{itemize}
        \item A foundation for enterprise \alert{Blockchain-as-a-Service (BaaS)} on Kubernetes.
      \end{itemize}
    \end{column}
  \end{columns}
  \vfill
  \textbf{Keywords:} \alert{Blockchain in Kubernetes} \;$\cdot$\; \alert{Enterprise Scale Blockchain} \;$\cdot$\; \alert{Private Blockchain}.
\end{frame}

% ===========================================================================
\begin{frame}{Keywords Recap}
  \small
  \begin{tabular}{@{}p{0.40\textwidth}p{0.54\textwidth}@{}}
    \toprule
    \textbf{Video} & \textbf{Keywords} \\
    \midrule
    What is Polkadot? (Beginners Guide) & Internet; Definition of Block; Decentralized System \\
    \addlinespace
    Managing \& Securing Blockchain on Kubernetes & Blockchain in Kubernetes; Enterprise-Scale Blockchain; Private Blockchain \\
    \bottomrule
  \end{tabular}
  \vfill
  \begin{itemize}
    \item Video 1: conceptual foundation --- what a block is and why decentralisation matters.
    \item Video 2: putting it into production --- permissioned blockchain, automated on Kubernetes, secured by isolation.
  \end{itemize}
\end{frame}

\begin{frame}[standout]
  Thank you! \\[0.5em]
  \small Questions?
\end{frame}

\end{document}
